\section{Exact target--synthesis spectrum}
\label{sec:tpc38-target-spectrum}

Fixed--target Parseval concerns the norm of one row.  A target--coupled
Bessel estimate concerns coherent combinations of many such rows.  The
distinction can be quantified without an inequality.

\subsection{The local full--face Gram}

Fix one odd prime \(q\), put \(n=q-1\), \(p=n^{-1}\), and let
\begin{equation}
 c_N(F)=\one_{F=0}-p\one_{F\ne-N},
 \qquad N\in\F_q^\times,\quad F\in\F_q.
 \label{eq:tpc38-local-full-row}
\end{equation}

\begin{lemma}[Exact local target Gram]
\label{lem:tpc38-local-target-Gram}
For \(N,N'\in\F_q^\times\),
\begin{equation}
 \langle c_N,c_{N'}\rangle
 =\begin{cases}
  1-p,&N=N',\\
  1-p-p^2,&N\ne N'.
 \end{cases}
 \label{eq:tpc38-local-target-Gram}
\end{equation}
Equivalently, the \(n\times n\) Gram matrix is
\begin{equation}
 G_q=p^2I_n+(1-p-p^2)\mathbf1\mathbf1^*.
 \label{eq:tpc38-local-Gram-matrix}
\end{equation}
Its squared singular values are
\begin{equation}
 \lambda_q(\eta)=
 \begin{cases}
  \displaystyle\frac{q(q-2)^2}{(q-1)^2},
  &\eta=\mathbf1,\\[0.9em]
  \displaystyle\frac1{(q-1)^2},
  &\eta\ne\mathbf1,
 \end{cases}
 \label{eq:tpc38-local-singular-values}
\end{equation}
where \(\eta\) ranges over target Mellin modes.
\end{lemma}

\begin{proof}
For \(N=N'\), the distinguished value \(F=0\) contributes
\((1-p)^2\), the killed value \(F=-N\) contributes zero, and the remaining
\(q-2=n-1\) values contribute \(p^2\).  Their sum is \(1-p\).

For \(N\ne N'\), the values \(0,-N,-N'\) are distinct.  The first
contributes \((1-p)^2\), the next two contribute zero, and the remaining
\(q-3=n-2\) values contribute \(p^2\).  The total is \(1-p-p^2\).
This proves \eqref{eq:tpc38-local-Gram-matrix}.  On the orthogonal
complement of \(\mathbf1\), its eigenvalue is \(p^2\).  On \(\mathbf1\),
the eigenvalue is
\[
 p^2+n(1-p-p^2)
 =\frac{q(q-2)^2}{(q-1)^2}.
\]
\end{proof}

The off--diagonal entry in \eqref{eq:tpc38-local-target-Gram} differs from
the row norm by only \(p^2\).  The local full rows are nearly parallel.

\subsection{The CRT tensor spectrum}

Return to three primes.  Define the full target--synthesis operator
\begin{equation}
 T_C:\ell^2((\Z/M\Z)^\times)\longrightarrow\ell^2(\Z/M\Z),
 \qquad
 (T_Ca)(F)=\sum_{N\bmod M}^{*}a_NC_N(F).
 \label{eq:tpc38-TC-definition}
\end{equation}
For a target Mellin mode
\(\eta=\eta_1\otimes\eta_2\otimes\eta_3\), put
\begin{equation}
 \lambda_C(\eta)=\prod_{i=1}^3\lambda_{q_i}(\eta_i),
 \label{eq:tpc38-lambda-C}
\end{equation}
with the local factors in \eqref{eq:tpc38-local-singular-values}.

\begin{theorem}[Complete full--face target spectrum]
\label{thm:tpc38-full-target-spectrum}
The target Mellin modes diagonalize \(T_C^*T_C\), and their squared
singular values are exactly \(\lambda_C(\eta)\).  In particular,
\begin{equation}
 \boxed{
 \|T_C\|_{\op}^2
 =\Lambda_C
 :=\prod_{i=1}^3
 \frac{q_i(q_i-2)^2}{(q_i-1)^2}.}
 \label{eq:tpc38-TC-norm}
\end{equation}
If \(q_i\asymp J\), then
\begin{equation}
 \Lambda_C\asymp\Phi\asymp J^3.
 \label{eq:tpc38-TC-endpoint}
\end{equation}
Under the same hypothesis \(q_i\asymp J\), every nonprincipal global
target mode has squared singular value \(O(1)\).
\end{theorem}

\begin{proof}
The CRT identifies \(T_C\) with the tensor product of the three local
matrices in \eqref{eq:tpc38-local-full-row}.  Tensor
\cref{lem:tpc38-local-target-Gram}.  The largest local eigenvalue is the
principal one.  If a global mode has a nonprincipal component at \(i\),
then that local factor is \(O(J^{-2})\), while the two remaining principal
factors are \(O(J)\).  Their product is \(O(1)\).
\end{proof}

The direct Gram formula is also useful.  For \(p_i=(q_i-1)^{-1}\),
\begin{equation}
 \boxed{
 \langle C_N,C_{N'}\rangle
 =\prod_{i=1}^3
 \left[
 (1-p_i)\one_{N=N'\ (\bmod q_i)}
 +(1-p_i-p_i^2)\one_{N\ne N'\ (\bmod q_i)}
 \right].}
 \label{eq:tpc38-global-full-row-Gram}
\end{equation}
If \(N\) and \(N'\) differ in every coordinate, the right side is
\(1-O(J^{-1})\).

\subsection{The proper target spectrum}

Define
\begin{equation}
 (T_Pa)(F)=\sum_{N\bmod M}^{*}a_NP_N(F).
 \label{eq:tpc38-TP-definition}
\end{equation}
Since \(P_N=\delta_0-C_N\) and \(C_N(0)=A\), the target Gram is
\begin{equation}
 T_P^*T_P=T_C^*T_C+(1-2A)\mathbf1\mathbf1^*.
 \label{eq:tpc38-proper-Gram-rank-one}
\end{equation}

\begin{theorem}[Complete proper--face target spectrum]
\label{thm:tpc38-proper-target-spectrum}
Every nonprincipal global target Mellin mode has squared singular value
\(\lambda_C(\eta)\).  The global principal mode has squared singular value
\begin{equation}
 \boxed{
 \lambda_P(\mathbf1)
 =\Lambda_C+\Phi(1-2A).}
 \label{eq:tpc38-proper-principal-eigenvalue}
\end{equation}
For \(q_i\asymp J\),
\begin{equation}
 \lambda_P(\mathbf1)\asymp\frac\Phi J\asymp J^2,
 \qquad
 \|T_P\|_{\op}^2\asymp J^2.
 \label{eq:tpc38-TP-endpoint}
\end{equation}
Moreover,
\begin{equation}
 \langle P_N,P_{N'}\rangle
 =1-2A+\langle C_N,C_{N'}\rangle,
 \label{eq:tpc38-proper-row-Gram}
\end{equation}
and
\begin{equation}
 \langle C_N,P_{N'}\rangle
 =A-\langle C_N,C_{N'}\rangle.
 \label{eq:tpc38-cross-row-Gram}
\end{equation}
\end{theorem}

\begin{proof}
Equation \eqref{eq:tpc38-proper-Gram-rank-one} follows by expanding
\(\langle\delta_0-C_N,\delta_0-C_{N'}\rangle\).  The rank--one correction
vanishes on every nonprincipal target mode and adds \(\Phi(1-2A)\) on the
principal mode.  This proves the exact spectrum.

For the asymptotic, divide \eqref{eq:tpc38-proper-principal-eigenvalue} by
\(\Phi\).  One has
\begin{equation}
 \frac{\Lambda_C}{\Phi}
 =\prod_{i=1}^3
 \left(1-p_i-p_i^2+p_i^3\right),
 \label{eq:tpc38-Lambda-density}
\end{equation}
while \(A=\prod_i(1-p_i)\).  Expanding the fixed products shows
\begin{equation}
 \frac{\lambda_P(\mathbf1)}\Phi
 =\sum_{i=1}^3p_i+O(J^{-2})\asymp J^{-1}.
 \label{eq:tpc38-proper-principal-density}
\end{equation}
The row Gram identities follow from the same expansion.
\end{proof}

\subsection{Trace is not synthesis}

The fixed--target row norms are
\begin{equation}
 \|C_N\|_2^2=A=1-O(J^{-1}),
 \qquad
 \|P_N\|_2^2=1-A\asymp J^{-1}.
 \label{eq:tpc38-row-norms-versus-synthesis}
\end{equation}
By contrast, the synthesis norms are of orders \(J^3\) and \(J^2\).
Thus a diagonal or trace estimate over \(N\) cannot be upgraded to a
target--coupled operator estimate without controlling the principal target
direction.  This is a precise finite--dimensional obstruction; it does not
assert that the actual M\"obius row vector occupies that direction.

\begin{corollary}[Coefficient--blind target Bessel obstruction]
\label{cor:tpc38-target-Bessel-obstruction}
Any coefficient--blind inequalities
\begin{equation}
 \|T_Ca\|_2^2\le B_C\|a\|_2^2,
 \qquad
 \|T_Pa\|_2^2\le B_P\|a\|_2^2
 \label{eq:tpc38-coefficient-blind-Bessel-bounds}
\end{equation}
valid for every scalar target vector \(a\) must satisfy
\begin{equation}
 B_C\ge\Lambda_C\asymp J^3,
 \qquad
 B_P\ge\lambda_P(\mathbf1)\asymp J^2.
 \label{eq:tpc38-coefficient-blind-lower-bounds}
\end{equation}
Thus the fixed--row norms alone cannot furnish smaller synthesis constants.
\end{corollary}

\begin{remark}[Scope]
The theorem concerns unrestricted scalar target coefficients on one CRT
period.  Physical coefficients are far more structured.  The result rules
out a coefficient--blind frame argument; it is neither a lower bound for
\(V_L+V_R\) nor evidence that the actual M\"obius vector is extremal.
\end{remark}
